%-------------------------------------------------------------------------------
%	SECTION TITLE
%-------------------------------------------------------------------------------
\cvsection{教育\! 背景}


%-------------------------------------------------------------------------------
%	CONTENT
%-------------------------------------------------------------------------------
\begin{cventries}

%---------------------------------------------------------
  \cventry
    {计算机科学 \textbar\ 硕士} % Degree
    {悉尼大学} % Institution
    {澳大利亚，新南威尔士州，悉尼} % Location
    {2024 年 2 月 - 2026 年 5 月} % Date(s)
    {
      \begin{cvitems} % Description(s) bullet points
        \item {\textbf{ISSRE 2026 录用论文（第二作者）：}\href{https://rahul.gopinath.org/publications/2026/07/11/drdd/}{``Dr. DD: 1-Minimal Isolation of Failure Causes via Deferred Restarts.''}}
        \item {\textbf{论文相关实现与实验：}实现\href{https://github.com/FreeFlyingSheep/delta-debugging}{Python Delta Debugging 缺陷输入归约框架}，自动缩减触发二进制软件故障的输入，辅助故障复现与定位；在 Binutils/Valgrind 上完成实验验证，对比归约效果、测试调用数与耗时。}
      \end{cvitems}
    }

  \cventry
    {计算机科学与技术（嵌入式培养）\textbar\ 本科} % Degree
    {苏州工学院（原常熟理工学院）} % Institution
    {中国，江苏省，苏州市} % Location
    {2016 年 9 月 - 2020 年 6 月} % Date(s)
    {
      \begin{cvitems} % Description(s) bullet points
        \item {\textbf{ICAIS 2020 论文（第一作者）：}\href{https://doi.org/10.1007/978-981-15-8083-3_25}{``Formal Verification of the Correctness of Operating System on Assembly Layer in Isabelle/HOL.''}}
        \item {\textbf{研究内容：}使用 Isabelle/HOL 对 MINIX 3 的 i386 上下文保存汇编代码建模，验证寄存器、内存状态及指令执行结果符合预期。}
      \end{cvitems}
    }

%---------------------------------------------------------
\end{cventries}
